% -----------------------------------------------------------------------------
% 第1章 序論（杉谷先生面談後半の論証順序に基づく再構成版）
% -----------------------------------------------------------------------------

\chapter{序論}
\label{chap:introduction}

\section{政策決定の合理化が残した問題}
\label{sec:intro_rationalization_problem}

\subsection{PPBSからEBPMまで}
\label{subsec:intro_ppbs_to_ebpm}

公共政策に計算機を導入すれば、より合理的な意思決定が可能になるという期待は、生成AIとともに初めて現れたものではない。1960年代の米国で展開された計画・プログラム・予算システム（Planning--Programming--Budgeting System: PPBS）は、目的、代替案及び費用を体系的に比較し、計画と予算を結び付けようとした。Schick（1966）は、予算改革を統制志向、管理志向、計画志向という段階から捉え、PPBSを計画機能に重点を置く改革として論じた\cite{schick1966ppb}。これに対してWildavsky（1969）は、分析に必要な情報、専門人材及び組織上の支持が一様に存在するわけではなく、国防部門で形成された仕組みを国内政策へ広く適用することの困難を指摘した\cite{wildavsky1969ppbs}。ここで示されたのは、分析手法を整備するだけでは、異なる目的を持つ組織を横断して情報を集め、その分析を実際の決定へ結び付ける条件まで整わないという問題である。

1980年代以降に広がった新公共管理（New Public Management: NPM）は、組織内の計画技法よりも、業績測定、成果管理、競争及び管理者の裁量によって公共サービスを改善しようとした。Hood（1991）はNPMの教義を整理するとともに、効率、公平、信頼性などの行政価値の間には競合があり、単一の管理原理をあらゆる公共サービスへ適用できるわけではないと論じた\cite{hood1991npm}。NPMは、投入と手続だけでなく成果を問う契機をもたらしたが、何を成果として測定するか、測定されにくい価値をどう扱うかという選択を不要にはしなかった。

2000年代以降に制度化されたエビデンスに基づく政策形成（Evidence-Based Policy Making: EBPM）は、政策手段の効果に関する知識を意思決定へ接続しようとする。Sanderson（2002）は、政策形成におけるエビデンス利用を道具的合理性だけに還元せず、実践的推論と社会的学習を含む過程として捉える必要を示した\cite{sanderson2002evidence}。Parkhurst（2017）も、科学的妥当性を欠くエビデンス利用を技術的バイアス、エビデンスを提示しやすい争点が他の社会的関心を周辺化することをイシュー・バイアスとして区別し、必要なのはエビデンスの増加だけでなく、その選択と利用を統治することだと論じる\cite{parkhurst2017politics}。EBPMが知識の質と利用を改善しても、どの問題を対象とし、どの価値に照らして結果を評価するかは、なお政策判断として残る。

PPBS、NPM及びEBPMは、対象、方法及び歴史的文脈を異にしており、一つの改革系列として同一視することはできない。それでも三つの蓄積をたどると、情報、測定又は分析能力の改善から、政策目的の選択、組織間の調整及び決定への責任が自動的に導かれるわけではないという問題が繰り返し現れる。本研究はこれを、合理化の失敗を宣言する根拠ではなく、計算機が政策形成のどこに作用し、どこに人間と社会の判断が必要となるかを特定する出発点とする。

政策研究は、包括的な分析によって一つの主体が政策全体を調整する構想とは異なる政策形成機構も論じてきた。Lindblom（1965）は、限定された知識と異なる利害を持つ複数主体が、互いの決定に反応しながら調整を重ねる partisan mutual adjustment を、民主的政策形成における知的機構として提示した\cite{lindblom1965intelligence}。足立（1998）はこれを「自己中心的相互調整」と訳し、反復的な政治過程における相互監視と要求調整の機構として整理している\cite{adachi1998inefficiency}。この議論は、政策形成に必要な知識を単一主体へ集約することだけでなく、分散した主体が互いの判断を観察し、異議を示し、修正できる環境を整える必要を示す。

\subsection{全面的な委任と全面的な排除の間}
\label{subsec:intro_between_delegation_and_rejection}

以上の政策合理化の経験と相互調整論から、アルゴリズムによる政策決定を一括して退ける結論は導かれない。PPBSへの批判も、政策分析そのものを不要としたのではなく、分析を有効にする組織条件と利用者の必要を問い直したものであった\cite{wildavsky1969ppbs}。NPMとEBPMへの批判も、成果の測定又はエビデンスの利用を放棄するのではなく、複数の行政価値と政治的関心を扱える制度を求めている\cite{hood1991npm,parkhurst2017politics}。政策研究が明らかにしてきたのは、計算を用いることの不可能性ではなく、計算によって政策判断の全体を置き換えることの困難である。

他方、今日のアルゴリズムを、PPBSが用いた分析技法又は事前に規則を記述するプログラムと同一視することもできない。生成AIは、数値化された変数だけでなく、法令、研究論文、行政文書、会議録及び当事者の語りを処理し、問題の分類、論点の比較、評価の記述及び選択肢の生成に利用できる。人間中心AIは、高い自動化と人間による統制を対立させず、信頼性、安全性及び人間の統制可能性を設計要件として組み込む方向を提示してきた\cite{shneiderman2022}。また、AIを用いた政策分析をめぐる研究は、エビデンス利用の拡張とともに、公平性、透明性、説明責任及び専門職倫理を扱う必要を指摘している\cite{newman2023mapping}。設計時に目的、原則及び制約を計算過程へ反映する余地が広がったことによって、何を計算機に担わせうるかは、過去の技術を前提とした議論だけでは決められなくなった。

したがって、本研究の課題は、アルゴリズムが政策決定を代替できるかという二者択一に答えることではない。計算機によって改善できる政策形成の機能、計算機の出力からは正当化できない判断、両者を接続する制度条件を区別することである。この区別には、計算機を一回の決定を補助する道具としてだけでなく、その生成物が文書、制度及び人間の判断へ取り込まれ、後続する判断条件を作り替える環境の構成要素として捉える視角が必要となる。そこで次に検討するのが、斉藤賢爾のメタ・ネイチャーである。

\section{メタ・ネイチャーが提起する問題}
\label{sec:intro_metanature_problem}

\subsection{「鉛筆が木に生える」人工環境}
\label{subsec:intro_pencil_tree}

鉛筆を一本作るにも、木材、黒鉛、塗料、金属、輸送及び販売に携わる多数の人々の分業が必要となる。斉藤賢爾は、この分業が極度に自動化され、鉛筆を木の実のように人工環境から得られる状態を構想した。人間が個々の生産工程を指示しなくても、人工環境が当初想定されなかった変化に対応しながら、生産と分配を継続する。この状態を斉藤はメタ・ネイチャー（Meta-Nature）と呼ぶ\cite[pp. 361--363]{saito2025metanature}。

この構想の魅力は、既存の作業を高速化することだけにない。現在の社会では、何かを実現するために、必要な専門知を持つ人や組織を探し、仕事を分配し、その成果を調整しなければならない。人工環境が新しい状況を認識し、必要な知識と機能を組み合わせ、生産と分配を続けられるならば、人間は希少な専門能力の所在によって活動可能性を制限されにくくなる。斉藤がメタ・ネイチャーによって描くのは、機械が人間の仕事を代替する未来にとどまらず、人間が人工環境との関係を通じて、何を行うべきかを構想し直せる社会である\cite[pp. 363--370]{saito2025metanature}。

公共政策もまた、専門知、法的権限、財源及び組織能力を分業によって組み合わせる営みである。自治体が交通計画を作成するときには、行政職員、交通事業者、住民、専門家、国及び広域自治体が、それぞれ異なる知識と権限を持ち寄る。人工環境が資料を検索し、状況を分類し、将来を予測し、政策案を生成し、その結果を次の判断へ反映するようになれば、変化するのは政策事務の速度だけではない。誰が何を問題として認識できるか、どの選択肢を構想できるか、何を評価基準として採用できるかという、政策形成の条件そのものが変わる。

本研究は、メタ・ネイチャーが公共政策の生産・分配及び形成過程に及ぶとき、政策を構想できる範囲に何が生じ、その変化を公共的に統治するために何を検討しなければならないかを問う。メタ・ネイチャーを政策研究へ適用する意義は、計算機を個別の政策手段として評価するだけでなく、計算機を含む人工環境の内部で、人間が政策目的を形成し、判断し、その環境を更新する関係を分析できる点にある。

\subsection{概念の系譜と近接概念}
\label{subsec:intro_metanature_lineage}

メタ・ネイチャーという語は、本研究の造語ではない。斉藤によれば、この概念は2019年のブロックチェーンに関する会議において、分業を前提とする鉛筆の生産を自動化によって組み替える着想から生まれた。その後、斉藤は2022年の国際会議論文で、人工環境から鉛筆を木の実のように得る構想を提示し\cite{saito2022path}, 2025年の論文では、人工環境が未知の状況に対応しながら生産と分配を自律的に継続する状態として概念を明示した\cite{saito2025metanature}。斉藤の2023年の論考は、デジタル技術を、固定した視点と均質な大量複製を強化してきた産業社会から、個々の主体による構成へ向かう反転として捉えており、メタ・ネイチャーを分業社会の変容として論じる思想的背景を示している\cite{saito2023reversegutenberg}。

本研究がメタ・ネイチャーを公共政策へ接続する技術的着想にも、ブロックチェーンとWeb3の影響がある。特定の中央管理者だけに依存せず、複数主体が規則、履歴及び状態を共有して検証する仕組みは、制度と記録を計算環境へ組み込む先例となった\cite{olnes2017blockchain,janssen2020blockchain}。この系譜から本研究が引き継ぐのは、既に決められた規則をコードで執行することではなく、人間、制度及び計算機の関係を、記録と検証を伴う環境として設計する発想である。この発想は、後に検討する秘密保護、判断理由の検証、異議申立て及び責任配分を、政策評価の制度と技術へ組み込む課題に接続する。\footnote{技術的系譜は、Bitcoin、スマートコントラクト、Ethereum及び分散型自律組織へとたどることができる。Nakamoto（2008）は、参加ノードの検証によって、信頼できる第三者なしに取引履歴の順序と改変困難性を共有するBitcoinを提示した\cite{nakamoto2008bitcoin}。Szabo（1996）が構想したスマートコントラクトは、契約条項を計算機化された取引プロトコルとして実行するものであり\cite{szabo1996smartcontracts}、Buterin（2014）のEthereumは、所有、取引、状態遷移及び組織運営の規則をスマートコントラクトとして記述できる基盤へ拡張した\cite{buterin2014ethereum}。その上では、提案、投票、委任、共同資金の配分及び規則変更を実装する分散型自律組織（Decentralized Autonomous Organization: DAO）が展開してきた\cite{hassan2021dao,hsieh2018dao}。}

斉藤の定義では、当初想定されなかった状況にも対応できる汎用人工知能（Artificial General Intelligence: AGI）が、メタ・ネイチャーの前提となる\cite[p. 361]{saito2025metanature}。したがって、現在の大規模言語モデルを利用しただけでメタ・ネイチャーが成立したことにはならない。本研究は、第1章でAGIの実現可能性又は性能上の限界を論証するのではなく、斉藤が示した人工環境が公共政策に及ぶ場合に生じる分析課題を提示する。生成AIが担いうる知識処理と、人間及び社会に帰属する判断との境界は、第5章で検討する。

メタ・ネイチャーが指す対象を明確にするには、近接する研究との違いを確認する必要がある。Kitchin and Dodge（2011）は、ソフトウェアが空間の外部に置かれた道具ではなく、日常生活の空間を継続的に作り替えることを示した\cite{kitchin2011codespace}。Yeung（2018）のアルゴリズム規制論は、環境から継続的にデータを収集して知識を生成し、所与の目標に向けてシステムの作動を修正する意思決定システムを分析する\cite{yeung2018algorithmicregulation}。Kephart and Chess（2003）の自律コンピューティングは、管理者が設定した目標に従って自らを管理する計算システムを構想した\cite{kephart2003autonomic}。Star and Ruhleder（1996）のインフラストラクチャー研究は、情報基盤が利用者、組織及び実践との関係の中で成立することを示している\cite{star1996infrastructure}。

これらの研究は、ソフトウェアによる空間の形成、継続的な計算による行為の調整、計算システムの自己管理及び基盤と人間の関係をそれぞれ説明する。これに対して斉藤のメタ・ネイチャーは、人工環境全体が人間にとって自然のような存在となり、生産と分配を継続する状態を対象とする。本研究がこの概念に注目するのは、公共政策に計算機を導入する場面だけでなく、計算機によって生成され続ける分類、予測、評価及び選択肢が、政策主体にとって所与の環境となる場面を捉えるためである。\footnote{落合（2018）のデジタルネイチャーは、計算機を人間の外部にある道具ではなく自然環境の一部として捉え、物理世界と情報世界の境界が失われる自然観を提示する\cite{ochiai2018digitalnature}。斉藤のメタ・ネイチャーは、人工環境が極度に自動化され、生産と分配を自律的に継続する作動状態に焦点を置く。本研究では、計算機が環境の一部となる点を両者の接点とし、生産・分配及び判断材料の継続的生成を扱う際にメタ・ネイチャーを用いる。}

\subsection{なぜ公共政策が扱う必要があるのか}
\label{subsec:intro_public_policy_relevance}

メタ・ネイチャーは、人工環境の将来像であると同時に、生産と分配の仕組みに関する社会構想である。交通、医療、福祉、防災、教育及び社会基盤の政策は、サービス、資源、危険及び機会を誰にどのように配分するかを決める。人工環境がこれらの生産と分配を継続的に担う場合、その目的、対象、優先順位及び例外処理は公共政策の問題となる。自動化された環境が安定して作動することと、その環境が行う配分が公共的に正当であることは同じではない。

さらに、生成AIは、既に決められた規則を実行するだけでなく、政策文書を要約し、問題を分類し、将来を予測し、評価を記述し、選択肢を生成する。Lee and Dai（2025）は、AIとデータ駆動技術が政策サイクルの各段階へ利用される一方、透明性、説明責任及び人間による監督が課題となることを示している\cite{lee2025ai_policy_cycle}。Janssen（2025）は、生成AIを、政策、データ、組織過程及び人間が相互に適応する複雑適応システムとして捉え、生成AIの利用に応じて統治の仕組みも更新する必要があると論じる\cite{janssen2025genaicas}。生成物が文書、データ及び制度運用へ取り込まれ、それらが将来の学習と判断の入力となるとき、計算機は政策形成の外部にある道具ではなく、後続する判断を条件付ける環境の一部となる。

この変化には、知識へのアクセスを広げる魅力がある。研究論文、行政文書、会議録及び当事者の語りを比較し、従来は担当者の経験又は専門家への接触可能性に左右されていた政策形成を支援できるからである。他方、出典と推論過程が不明確な生成物が計画及び評価へ反復して取り込まれれば、低品質で検証困難なAIスロップ（AI slop）が次の判断の前提として蓄積する\cite{cambridgeaislop2026}。生成AIは、政策形成に利用できる知識を増やすと同時に、何を信頼し、何を政策目的として選び、その結果に誰が責任を負うかという問題を拡大する。

したがって、公共政策研究に求められるのは、メタ・ネイチャーの到来を技術的に予測することだけではない。人工環境が政策主体の判断条件となる過程を観察し、そこに組み込まれる目的、知識、権限、参加、記録及び責任を検討することである。そのためには、人工環境の作用先を、公共政策研究の分析概念によって特定しなければならない。

\section{政策デザイン空間への作用}
\label{sec:intro_policy_design_space}

\subsection{政策デザイン空間とは何か}
\label{subsec:intro_design_space_definition}

公共政策では、目的が先に確定し、利用可能な手段だけが比較されるとは限らない。政策主体が何を問題と認識し、どの目的を追求し、どの手段を選択肢として構想できるかは、制度的遺産、政治関係、知識及び利用可能な資源によって方向付けられる。Howlett and Mukherjee（2014）が提示した政策デザイン空間（policy design space）は、政府が政策を意図的に設計しようとする程度と、設計に利用できる知識・資源の組合せから、政策形成が意図的なデザインと非デザインのどこに位置するかを説明する\cite{howlett2014policydesign}。Howlettらはさらに、政策形成環境、政策ミックス及び個別政策手段という複数の水準から、政策デザインが成立する条件を整理した\cite{howlett2018designspaces}。では、計算機を含む人工環境は、この空間のどこに作用するのか。

計算機が同じ資料をより速く処理するだけなら、その変化は既存の政策デザイン空間の内部における分析能力の向上として説明できる。しかし、計算機が参照可能な資料を選び、問題の分類を提示し、将来予測を反復し、評価指標と政策案を生成し、その生成物が制度と記録へ組み込まれる場合、政策主体が構想できる目的と選択肢の範囲自体が変わる。したがって、メタ・ネイチャーが公共政策に及ぼす作用を観察する中心的な対象は、政策デザイン空間である。問題は、その作用が政策主体の構想可能性を広げるのか、狭めるのかである。

人工環境が従来アクセスできなかった知識と選択肢を提示すれば、政策デザイン空間は拡張する。反対に、既存の指標、文書様式及び多数派の判断を反復すれば、その空間は狭窄する。本研究は、この両方向の作用と、それを公共的に統治する条件を検討する。両者を区別するには、選択肢の広狭だけでなく、その空間を成立させる能力、知識の更新、知識の流通及び主体間関係を分析しなければならない。

\subsection{政策デザイン空間を捉える四つの視角}
\label{subsec:intro_four_lenses}

この分析は、まず、政策主体がどの能力と資源を利用できるかを確かめることから始まる。政策能力、政府学習、EIPMエコシステム及びメタガバナンスは、同じ現象に別名を与える理論ではなく、能力の所在から知識の更新と流通を経て、主体間関係の統治へ進む四つの理論的視角である。

Wu, Ramesh and Howlett（2015）は、政策能力を分析的・運用的・政治的能力の三側面と、個人・組織・システムの三水準から整理した\cite{wu2015policycapacity}。個人の分析技能だけでなく、組織が利用できる情報資源、実施能力、信頼、調整及び政治的支持が政策形成を支える。Mukherjee, Coban and Bali（2021）は、この枠組みを政策デザイン空間へ接続し、有効な政策デザインが複数水準の能力に支えられることを示した\cite{mukherjee2021policycapacities}。政策能力論によって、人工環境が能力を補完するのか、特定の組織へ能力を集中させるのかを検討できる。しかし、能力の所在を確認するだけでは、政策に用いる知識がどのように更新されるかは分からない。

この更新過程を扱うのが政府学習論である。唱道連合フレームワークは、政策志向学習を、信念体系に関わる認識の持続的変化として政策変化の内部に位置付けた\cite{sabatier1988}。秋吉（2021）は、政府学習を情報の獲得だけでなく、政策知識の共有・格納及び政策案形成からなる過程として再構成した\cite{akiyoshi2021governmentlearning}。人工環境が生成した要約、評価及び選択肢が組織的記憶へ蓄積する場合、政府学習論によって、何が学習され、何が失われ、後続する政策案がどの知識に依存するかを分析できる。それでも、組織内部の学習だけを見れば、誰が知識を作り、仲介し、利用可能にするかという流通関係が捉えられない。

知識の流通関係を捉えるために、EIPMエコシステムの視角が必要となる。OECDと欧州委員会（2025）は、個人、組織、組織間及びシステムの各水準を区別し、共有言語、継続的能力形成、知識仲介及び多様なエビデンスへのアクセスを、水準を横断する課題として示した\cite{oecd2025eipmecosystems}。生成AIが検索と要約を担っても、異なる資料の作成目的、検証方法、公開可能性及び時間軸の違いは残る。EIPMエコシステムによって、人工環境が誰の知識を接続し、誰の経験を不可視化するかを検討できる。他方、知識の需要、供給及び仲介を記述しても、参加規則と知識利用の責任を誰が定めるかという問題は残る。

この問題に応答するのがメタガバナンス論である。S{\o}rensen and Torfing（2009）は、ネットワークの形成、相互作用のルール、資源配分及び参加者間関係を調整するメタガバナンスによって、ガバナンスネットワークの有効性と民主性を両立させる条件を論じた\cite{sorensen2009metagovernance}。人工環境が複数主体へ政策案と評価を提示する場合、参加規則、異議申立て、監査及び責任配分を定めなければ、その出力が事実上の決定となる。メタガバナンス論によって、人間と計算機を含む政策形成過程を公共的に統治する制度条件を検討できる。

\begin{table}[H]
  \centering
  \caption{政策デザイン空間を捉える四つの理論的視角}
  \label{tab:four_lenses}
  \small
  \begin{tabularx}{\textwidth}{p{0.21\textwidth}p{0.23\textwidth}X}
    \toprule
    理論的視角 & 説明する局面 & その視角だけでは残る問題 \\
    \midrule
    政策能力 & 能力と資源の所在、補完及び集中 & 知識が更新される過程 \\
    政府学習 & 知識の獲得、共有、格納及び更新 & 知識を作り、仲介し、利用可能にする関係 \\
    EIPMエコシステム & 知識の需要、供給、仲介及びアクセス & 参加規則と知識利用の責任 \\
    メタガバナンス & 主体間関係、参加規則及び責任の調整 & 四つの局面と人工環境の相互関係 \\
    \bottomrule
  \end{tabularx}
\end{table}

表\ref{tab:four_lenses}が示すように、四つの視角は、前の視角が説明した局面を引き継ぎながら、その視角だけでは残る問題を次の分析対象へ変える。この関係を用いることで、人工環境が政策デザイン空間へ及ぼす作用を、技術性能、組織能力又は知識利用のいずれか一つへ還元せずに分析できる。

\subsection{人間と人工環境による継続的な更新}
\label{subsec:intro_recursive_update}

政策主体は、地形、人口、制度、財政、人間関係、知識及び記録から判断材料と制約を受け取る。同時に、その判断は政策、制度、予算、評価及び文書として残り、後続する政策主体の判断条件となる。計算機がこの関係へ組み込まれると、既存の資料と人間の判断を入力として、分類、予測、評価及び選択肢を生成し、その生成物が次の判断材料として環境へ戻る。計算機が一度だけ答えを提示するのではなく、人間の判断を取り込みながら生成を続ける点に、政策デザイン空間の継続的な更新が生じる。

\begin{figure}[H]
  \centering
  \includegraphics[width=0.98\textwidth]{ch1/metanature_policy_design_space.png}
  \caption{メタ・ネイチャーと政策デザイン空間の相互関係（著者作成）}
  \label{fig:metanature_policy_design_space}
\end{figure}

図\ref{fig:metanature_policy_design_space}の中心には、政策主体が問題、目的、評価基準及び選択肢を構想する政策デザイン空間がある。制度・政治、知識・記録、資源・能力及び主体間関係がその範囲を形成し、計算機を含む人工環境が分類、予測、評価及び選択肢を生成する。政策主体の判断は人工環境へ記録され、その後の生成を変え、更新された生成物が次の政策デザイン空間を形成する。この相互関係が、生産・分配を継続する人工環境へ広がるとき、斉藤が示したメタ・ネイチャーが公共政策に及ぶ局面として観察できる。

ここで公共政策上の論点となるのは、計算機が存在するか否かではない。政策デザイン空間を広げる知識接続、目的と評価基準を選び直す規範的判断、認知的制約を持つ主体間の協調、秘密保護と検証可能性、異議申立てと責任配分を、人工環境の作動へどのように組み込むかである。次節では、この問題を診断する事例として地域公共交通政策を取り上げる。

\section{地域公共交通政策による診断}
\label{sec:intro_transport_case}

\subsection{権限移譲と判断条件}
\label{subsec:intro_decentralization}

日本の中央・地方関係は、政策デザイン空間が法的権限だけでは決まらないことを示す。2000年に施行された地方分権一括法は、機関委任事務制度を廃止し、国と地方自治体の関係を「上下・主従」から「対等・協力」へ転換した\cite{chihobunken1999}。この改革によって自治体が政策を判断する法的空間は拡大した。他方、2024年の地方自治法改正は、大規模災害、感染症のまん延その他これらに類する事態について、個別法に根拠がない場合にも国が自治体へ必要な措置を指示できる「国の補充的な指示」を設けた\cite{chihojichiho2024,chihoseido33,sangiin2024localgov}。

この改正から、自治体が能力を獲得しなかったために国の介入が必要になったとは結論できない。改正の直接的な背景には、コロナ禍における情報共有、役割分担及び責任所在の不明確さがあった\cite{chihoseido33,sangiin2024localgov}。この事例が示すのは、権限の所在を決めても、必要な知識、組織間関係及び非常時の調整方法が同時に整うわけではないことである。政策デザイン空間を分析するには、誰が権限を持つかに加えて、その主体が何を目的として構想でき、どの条件で他の主体が判断を補充又は変更できるかを検討しなければならない。

平時の計画行政にも、権限と判断条件のずれが現れる。内閣府調査によれば、2022年末時点で法令に基づく自治体計画の策定規定は524条項あり、その内訳は義務206条項、努力義務90条項、策定可能とする規定229条項であった。策定可能規定の約7割、努力義務規定の約4分の1は、補助、資金、税制又は規制特例等の要件と結び付いている\cite{cao2023localplans}。任意又は努力義務という法形式であっても、財政措置へ接続するには、国が定めた計画、記載事項及び評価方法を受け入れる必要が生じる。

政府は、総合計画や他の個別計画との一体策定を認め、計画策定の負担軽減を進めてきた。2024年の見直し後、他計画との一体策定が可能な計画は都道府県で約8割、市町村で約7割に増えた\cite{cao2023localplans,cao2024integratedplans}。しかし、市区町村の15.1％は一体策定の方法が分からず、22.9％は個別に策定するより負担になると回答した\cite{cao2024municipalplans}。計画期間、様式、所管部局、審議会手続及び補助要件が一致しない場合、複数計画を一つの文書にまとめても、政策目的と判断過程は統合されない。

\subsection{地域公共交通政策を事例とする理由}
\label{subsec:intro_case_fit}

本研究は、人工環境が作用する以前の政策デザイン空間を診断し、その後に人間と計算機の役割及び制度・技術設計を検討するため、地域公共交通政策を主要な事例とする。この政策領域では、移動の効率性、採算性、交通弱者の生活機会、地域拠点の維持、脱炭素及び交通事業者の持続可能性が競合する。単一の外部指標から政策目的を決めることは難しく、行政の計画と事業記録、交通工学・経済学の研究、住民・利用者の経験、事業者の運行・営業情報を比較しなければならない。

地域公共交通は、行政、住民、事業者、専門家、国及び広域自治体が相互依存する政策領域でもある。自治体だけでは路線を運行できず、事業者だけでは公共的な移動保障を決められない。住民参加だけでも、運行上の制約と財政負担を処理できない。しかも、事業者が保有する原価、需要及び運行上の情報には公開できない部分があるため、必要な知識の利用と秘密保護を同時に扱う必要がある。政策能力、政府学習、知識の需要・供給・仲介及び主体間関係が、一つの政策過程に現れる事例である。

地域公共交通政策では、国が自治体へ計画策定を求める範囲が拡大してきた。2007年制定の地域公共交通活性化再生法は、市町村が地域公共交通総合連携計画を作成できる制度を設けた。2014年改正は、同計画を地域公共交通網形成計画へ改め、都道府県を策定主体に加えた\cite{mlit2014networkplan}。2020年改正は名称を地域公共交通計画へ変更し、原則として全地方公共団体に作成の努力義務を課した。計画作成後には施策の実施状況を毎年度調査、分析及び評価する努力義務があり、利用者数、収支等の定量的目標も定めるよう努めることとされた\cite{regionaltransportact,mlit2020planobligation}。2023年改正は、地域関係者間の連携に関する事項を計画へ記載する努力義務を追加した\cite{mlit2023redesign}。

制度の拡張に伴い、地域公共交通計画の作成件数は2023年12月の901件から、2024年3月の1,021件、2024年度末の1,184件へ増加した\cite{mlit2024substantialization,mlit2025plantrend}。計画制度は予算執行とも接続している。地域公共交通確保維持改善事業は、地域公共交通計画へ位置付けられた補助対象事業について、実施状況と目標達成状況の評価を求める。地域公共交通利便増進実施計画の認定を受けた事業には、地域間幹線系統の補助要件の特例や、利用促進・利便性向上に関する支援が設けられている\cite{mlittransportsubsidy,mlitredesignbudget}。計画は政策方針を記述する文書であると同時に、個別事業を国の補助制度へ接続する入口となる。

地方バスに関する特別交付税措置にも、国の政策体系へ接続する誘因がある。山本（2023）によれば、自治体が不採算バス路線の維持に支出した経費は長く原則8割が算定対象とされてきたが、2021年の省令改正後、国の地域公共交通確保維持改善事業と連携しない事業は財政力指数に応じた減額算定の対象となった\cite{yamamoto2023busgrant,specialallocationtaxregulation}。財政力指数が高い自治体ほど、国庫補助体系と別に実施する事業への措置率が低くなる。計画策定の努力義務、認定事業への予算措置及び特別交付税の算定規則が重なることで、法的には自治体の選択に委ねられた政策が、財政面では国の制度メニューへ方向付けられる。

以上の特徴により、地域公共交通政策では、自治体が形式上の権限を持つことと、独自の問題、目的、評価基準及び選択肢を構想できることの違いを観察できる。第2章から第4章は、MaaSの導入、自治体調査及び全国の評価指標分析を通じて、現在の政策デザイン空間が制度、知識及び政府間関係によってどのように形成されているかを診断する。この診断によって、人工環境が政策形成へ深く組み込まれる前に、補完すべき能力と知識接続、及び固定化を避けるべき目的、評価基準と主体間関係を特定する。

\section{リサーチクエスチョン}
\label{sec:intro_research_questions}

メタ・ネイチャーが公共政策に及ぶ場合、人工環境が政策デザイン空間へ作用すること自体を確認するだけでは足りない。第一に、現在の政策目的と評価基準が、どの制度、知識及び組織間関係の下で形成されているかを診断する必要がある。診断しなければ、人工環境が何を補い、何を反復しているかを区別できないからである。

第二に、生成AIが支援できる知識処理と、人間及び社会的手続が担う判断を区別する必要がある。生成AIは資料を比較して選択肢を提示できるが、その処理能力だけから、どの価値を優先する権限、選択を正当化する責任及び異議を受ける立場は導かれない。この理論的境界を検討するのが第5章である。

第三に、判断を人間へ帰属させても、集合的な政策判断が直ちに成立するわけではない。政策主体は、確証バイアス、現状維持バイアス及び限られた視野を持ち、互いに異なる利害と知識を持つ。人工環境が複数の主体を接続するためには、認知的制約を持つ主体の協調がどの条件で維持されるかを検討しなければならない。この関係を限定されたモデル条件の下で分析するのが第6章である。

第四に、診断、役割境界及び協調条件を、実際の制度と技術へ変換する必要がある。政策評価では、事業者の秘密情報を利用しながら、評価理由を検証し、市民が評価原則へ関与し、異議を申し立てられる仕組みが求められる。秘密保護、検証可能性及び民主的正当性を同時に扱う政策評価環境を設計し、その作動を検証するのが第7章と第8章である。

以上の論証を、本研究は次の四つのリサーチクエスチョン（Research Question: RQ）として構成する。

\begin{table}[H]
  \centering
  \caption{本研究のリサーチクエスチョン}
  \label{tab:rq_general}
  \small
  \begin{tabularx}{\textwidth}{p{0.10\textwidth}X}
    \toprule
    RQ & リサーチクエスチョン \\
    \midrule
    RQ1 & 政策目的と評価基準の形成は、どのような制度、知識及び組織間関係によって制約されるか \\
    RQ2 & 生成AIは政策形成のどの機能を支援でき、どの規範的判断を人間と社会に残すべきか \\
    RQ3 & 認知的制約を持つ複数主体の協調は、どの条件で維持されるか \\
    RQ4 & 秘密保護、検証可能性及び民主的正当性を同時に扱う政策評価環境を設計できるか \\
    \bottomrule
  \end{tabularx}
\end{table}

RQ1が現在の政策デザイン空間を診断し、RQ2が人工環境におけるAIと人間の役割を区別し、RQ3が人間による協調の成立条件を検討し、RQ4がその知見を制度と技術へ移す。この順序によって、政策目的、参加、正当性及び責任を備えた人工環境をどのように統治できるかを検討する。

\section{本研究の位置づけと論証構造}
\label{sec:intro_position}

\subsection{先行研究に対する位置づけ}
\label{subsec:intro_contribution}

本研究は、斉藤が生産・分配及びイノベーションの将来像として提示したメタ・ネイチャーの定義に従い、人工環境が公共サービスの生産・分配及び政策形成へ及ぶ局面を分析する。メタ・ネイチャーを公共政策へ接続する媒介となるのが政策デザイン空間であり、人工環境が問題認識、目的、評価基準及び選択肢の構想可能性をどのように変えるかを問う。

政策デザイン研究に対しては、計算機を新たな政策手段又は分析能力としてだけでなく、政策デザイン空間を継続的に更新する環境として位置付ける。政策能力、政府学習、EIPMエコシステム及びメタガバナンスの研究に対しては、それぞれの知見を一つの包括理論へ統合するのではなく、人工環境が能力、知識の更新、知識の流通及び主体間関係へ及ぼす作用を観察する視角として用いる。これにより、技術性能、組織能力、知識利用及び民主的統治を、政策デザイン空間の変化という共通の対象について検討する。

地域公共交通政策の研究に対しては、計画の有無又は個別施策の効果だけでなく、自治体が何を政策目的とし、何を評価基準として構想できるかを形成する制度、様式、知識及び政府間関係を明らかにする。生成AIと政策形成の研究に対しては、AIを最終判断者として置くのではなく、資料と経験の接続、論点と選択肢の生成、評価理由の比較及び記録を担う環境構成要素として位置付ける。さらに、その利用を秘密保護、検証、異議申立て及び責任配分を備えた制度・技術として試行する。

\subsection{論文全体の論証構造}
\label{subsec:intro_argument_structure}

本論文の論証は、現状の診断、理論的境界、協調条件、制度・技術設計及び統合という順序で進む。第2章から第4章は、地域公共交通政策を対象として現在の政策デザイン空間を診断する。第5章は、生成AIが担いうる知識処理と、人間及び社会に帰属する判断との境界を検討する。第6章は、認知的制約を持つ複数主体の協調条件を計算論的に分析する。第7章と第8章は、これらの知見を政策評価の制度と技術へ変換し、Verdictによる実験室内実証及び実地実験の設計を行う。第9章は各RQへの応答を統合し、人工環境が政策デザイン空間へ及ぶ場合の制度設計原則を導く。第10章は、研究全体への回答と一般化の範囲を示す。

\begin{table}[H]
  \centering
  \caption{RQ、方法、対象章及び論証上の役割}
  \label{tab:rq_structure}
  \small
  \begin{tabularx}{\textwidth}{p{0.08\textwidth}p{0.16\textwidth}p{0.24\textwidth}X}
    \toprule
    問い & 対象章 & 主な方法 & 論証上の役割 \\
    \midrule
    RQ1 & 第2--4章 & 事例分析、自治体調査、評価指標の定量分析 & 現在の政策デザイン空間を形成する制度、様式、知識流通及び政府間関係を診断する \\
    RQ2 & 第5章 & 理論分析 & 生成AIが支援できる知識処理と、人間・社会へ帰属する判断を区別する \\
    RQ3 & 第6章 & 計算論的モデル実験 & 限定されたモデル条件下で、認知的制約と複数主体の協調の関係を検討する \\
    RQ4 & 第7--8章 & システム設計、実験室内実証、実地実験設計、秘密保護の限定検証 & 知識接続、秘密保護、検証可能性及び民主的正当性を扱う環境の部分的な実現可能性を検討する \\
    \bottomrule
  \end{tabularx}
\end{table}

\subsection{各章の構成}
\label{subsec:intro_chapter_outline}

第2章は、日本版「サービスとしてのモビリティ」（Mobility as a Service: MaaS）の導入過程を対象に、交通サービスの技術的統合と交通まちづくりの政策目的との間に生じる乖離を分析する。第3章は、北海道・東北運輸局管内42自治体の調査から、地域公共交通政策の目的と費用便益比等の評価基準との関係を検討する。第4章は、全国987自治体の12,568指標を分析し、評価指標の選択にみられる運輸局間の差異と制度的同型化を検討する。三つの章は、自治体の政策デザイン空間が、地域特性だけでなく、制度、知識及び政府間関係によって形成される過程を異なる資料と方法から診断する。

第5章は、創造システム理論と社会システム理論を用いて、生成AIが担いうる知識処理と、人間及び社会に帰属する判断との境界を論じる。第6章は、協調ロボット制御モデルによる22,000回のシミュレーションから、確証バイアス、現状維持バイアス及び狭い視野が協調指標へ及ぼす影響を検討する。このモデルは現実の政策過程又は熟議の代理ではなく、認知的制約を持つ主体間協調のメカニズムを限定条件下で分析する。

第7章は、競合する事業者等の秘匿情報を保護しながら、判断原則、評価理由、検証可能性及び異議申立てを維持する政策評価システムを設計する。明示された原則集合に基づいてAIモデルを訓練・評価するConstitutional AIと、LLMを評価者として用いるLLM-as-a-Judgeは、最終判断者ではなく、原則に照らして資料と政策案を比較する機構として位置付ける。第8章はVerdictを対象として、盛岡市民アンケートに由来する実験室内実証と秘密保護の限定検証を行い、岩手県内の実地実験の手続を設計する。

第9章は、四つのRQへの応答を統合し、メタ・ネイチャーが公共政策に及ぶ場合の統治原則を導出する。第10章は研究全体を総括し、地域公共交通政策への適用、他の政策領域へ一般化できる範囲及び今後の検証課題を示す。

本研究が検証するのは、地域公共交通政策から診断した現在の政策デザイン空間を起点として、生成AIと人間の役割、主体間協調及び政策評価の制度・技術を順に接続できるかである。この検証から、人工環境が公共政策に及ぶ際に必要となる統治上の論点を明らかにする。これによって、メタ・ネイチャーを技術的な未来像にとどめず、人間が人工環境の内部で目的を形成し、判断し、異議を述べ、その環境を更新する公共政策の問題として提示する。

\printbibliography[heading=subbibliography,title={章末参考文献}]
